
  
% ------------------------------------------------------------------------
% ------------------------------------------------------------------------
% abnTeX2: Modelo de Projeto de pesquisa em conformidade com 
% ABNT NBR 15287:2011 Informação e documentação - Projeto de pesquisa -
% Apresentação 
% ------------------------------------------------------------------------ 
% ------------------------------------------------------------------------
 
\documentclass[
	% -- opções da classe memoir --
	12pt,				% tamanho da fonte
	openright,			% capítulos começam em pág ímpar (insere página vazia caso preciso)
	oneside,			% para impressão em verso e anverso. Oposto a oneside
	a4paper,			% tamanho do papel. 
	% -- opções da classe abntex2 --
	%chapter=TITLE,		% títulos de capítulos convertidos em letras maiúsculas
	%section=TITLE,		% títulos de seções convertidos em letras maiúsculas
	%subsection=TITLE,	% títulos de subseções convertidos em letras maiúsculas
	%subsubsection=TITLE,% títulos de subsubseções convertidos em letras maiúsculas
	% -- opções do pacote babel --
	english,			% idioma adicional para hifenização
	french,				% idioma adicional para hifenização
	spanish,			% idioma adicional para hifenização
	brazil,				% o último idioma é o principal do documento
	]{abntex2}

% ---
% PACOTES
% --- 

% ---
% Pacotes fundamentais 
% ---
\usepackage{cmap}
\usepackage{lmodern}			% Usa a fonte Latin Modern
\usepackage{graphicx,url}			% Inclusão de gráficos
\usepackage[utf8]{inputenc}
\usepackage[brazil]{babel}		% Codificacao do documento (conversão automática dos acentos)
\usepackage[T1]{fontenc}		% Selecao de codigos de fonte.
\usepackage{lastpage}
\usepackage{indentfirst}		% Indenta o primeiro parágrafo de cada seção.
\usepackage{color}				% Controle das cores
\usepackage{microtype} 			% para melhorias de justificação
\usepackage{subfigure}
\usepackage{placeins}
\usepackage{hyperref}		% controla a formação do índice
\usepackage[portugues,ruled,lined]{algorithm2e}
\usepackage{algorithmic}
\usepackage{amsfonts, amssymb, amsmath, dsfont}		% fontes e símbolos
% matemáticos
\usepackage{mathtools}
\usepackage{verbatim}			% Permite apresentar texto tal como escrito no
% documento, ainda que sejam comandos Latex
\usepackage{multirow, array}	% Permite tabelas com múltiplas linhas e colunas
\usepackage{float}				% Necessário para tabelas/figuras em ambiente
% multi-colunas ---
\usepackage{bookmark}			% Cria menu de bookmarks

% Pacotes adicionais, usados apenas no âmbito do Modelo Canônico do abnteX2
% ---
%\usepackage{lipsum}				% para geração de dummy text
% ---
 
% ---
% Pacotes de citações
% ---
\usepackage[brazilian,hyperpageref]{backref}	 % Paginas com as citações na bibl
\usepackage[alf]{abntex2cite}	% Citações padrão ABNT
\usepackage{lscape}

% --- 
% CONFIGURAÇÕES DE PACOTES
% --- 
 
% ---
% Configurações do pacote backref
% Usado sem a opção hyperpageref de backref
%\renewcommand{\backrefpagesname}{Citado na(s) página(s):~}
%% Texto padrão antes do número das páginas
%\renewcommand{\backref}{}
%% Define os textos da citação
%\renewcommand*{\backrefalt}[4]{
%	\ifcase #1 %
%		Nenhuma citação no texto.%
%	\or
%		Citado na página #2.%
%	\else
%		Citado #1 vezes nas páginas #2.%
%	\fi}%
%% ---

% ---
% Informações de dados para CAPA e FOLHA DE ROSTO
% ---
\titulo{Detecção de Fraudes em Sistemas de Gestão Agropecuária: Uma Abordagem Baseada em Sistemas Inteligentes}
\autor{Leonardo Vieira Guimarães}
\local{Belo Horizonte}
\data{Janeiro de 2025}
\cefet{Centro Federal de Educação Tecnológica de Minas Gerais -- CEFET/MG}
\programa{Programa de Pós-graduação em Modelagem Matemática e Computacional}
\orientador{} % Orientador não foi especificado
%\coorientador{}
\instituicao{%
  \imprimircefet \par \imprimirprograma \par
}
\tipotrabalho{ANEXO III – Projeto de Pesquisa}
\tituloAcademico{Mestre} % Pode ser alterado para Doutor se necessário

% O preambulo deve conter o tipo do trabalho, o objetivo, o nome da instituição
% e a área de concentração
% Preâmbulo agora definido no arquivo de folha de rosto
%\preambulo{Projeto de Pesquisa apresentado no âmbito do Processo Público para
%         Seleção e Admissão de Novos Alunos Regulares para o Programa de Pós-Graduação em Modelagem 
%         Matemática e Computacional do Centro Federal de Educação Tecnológica de
%         Minas Gerais.}

% ---

\newtheorem{teo}{Teorema}[section]
\newtheorem{lema}[teo]{Lema}
\newtheorem{cor}[teo]{Corol\'ario}
\newtheorem{prop}[teo]{Proposi\c{c}\~ao}
\newtheorem{defi}[teo]{Defini\c{c}\~ao}
% ---
% Configurações de aparência do PDF final

% alterando o aspecto da cor azul
\definecolor{black}{RGB}{41,5,195}

% informações do PDF
\makeatletter
%\hypersetup{
		%pagebackref=true,
		%pdftitle={\@title}, 
	%	pdfauthor={\@author},
	%	pdfsubject={\imprimirpreambulo},
		%pdfcreator={LaTeX with abnTeX2},
	%	pdfkeywords={abnt}{latex}{abntex}{abntex2}{projeto de pesquisa}, 
		%colorlinks=true,       		% false: boxed links; true: colored links
		%linkcolor=black,          	% color of internal links
		%citecolor=black,        		% color of links to bibliography
		%filecolor=magenta,      		% color of file links
		%urlcolor=black,
		%bookmarksdepth=4
%}


\hypersetup{
	colorlinks,
	citecolor=black,
	filecolor=black,
	linkcolor=black,
	urlcolor=black
}

\makeatother
% --- 

% --- 
% Espaçamentos entre linhas e parágrafos 
% --- 

% O tamanho do parágrafo é dado por:
\setlength{\parindent}{1.3cm}

% Controle do espaçamento entre um parágrafo e outro:
\setlength{\parskip}{0.0cm}  % tente também \onelineskip

% ---
% compila o indice
% ---
\makeindex
% ---

% ----
% Início do documento
% ----
\begin{document}

% Seleciona o idioma do documento (conforme pacotes do babel)
%\selectlanguage{english}
\selectlanguage{brazil}

% Retira espaço extra obsoleto entre as frases.
\frenchspacing 

% ----------------------------------------------------------
% ELEMENTOS PRÉ-TEXTUAIS
% ----------------------------------------------------------
% \pretextual

% ---
% Folha de rosto (primeira página)
% ---
\include{./01-Elementos-pre-textuais/FolhaRosto/folhaRosto}
\imprimirfolhaderosto
% --- 





%\tableofcontents
%\chapter{Resumo}
\addcontentsline{toc}{chapter}{Resumo}
\label{chap:resumo}


\begin{resumo}

Este projeto propõe o desenvolvimento de algoritmos inovadores baseados em sistemas fuzzy evolutivos para detecção e prevenção de fraudes em sistemas de gerenciamento agropecuário. O diferencial dos sistemas fuzzy evolutivos está na capacidade de adaptar, de forma autônoma, tanto sua estrutura quanto seus parâmetros, à medida que novos dados são recebidos em fluxo contínuo. A modelagem fuzzy permite lidar com incertezas e ambiguidades presentes nos dados reais, utilizando funções de pertinência para representar graus de pertencimento a diferentes padrões ou comportamentos, o que é fundamental para identificar fraudes em cenários complexos e dinâmicos. A crescente ocorrência de fraudes fiscais nesse segmento, evidenciada por operações como CIRA, Rei do Gado e Boi na Linha, reforça a necessidade de abordagens tecnológicas inovadoras. O trabalho será desenvolvido em três etapas: (1) revisão da literatura sobre sistemas inteligentes adaptativos, detecção de anomalias e aprendizado incremental; (2) proposição, desenvolvimento e implementação de algoritmos fuzzy evolutivos, com foco na construção e atualização dinâmica das regras e funções de pertinência; (3) avaliação e validação dos algoritmos propostos em comparação com métodos do estado da arte, utilizando dados reais do setor agropecuário. Espera-se como resultado a obtenção de soluções com alta adaptabilidade, acurácia e interpretabilidade, além de contribuições teóricas e práticas para a detecção de fraudes, com potencial de aplicação em outros domínios de gestão pública.
 
 \vspace{0.5cm}
 
\textbf{Palavras-chave}: Detecção de fraudes em Sistemas, Sistemas Fuzzy Evolutivos, Pertinência, Setor agropecuário.

\end{resumo}
  
% --------------------------------------------------------------------------------------
% Sumario
% --------------------------------------------------------------------------------------
\include{./01-Elementos-pre-textuais/Sumario/sumario}

% ----------------------------------------------------------
% ELEMENTOS TEXTUAIS
% ----------------------------------------------------------
\textual
% ----------------------------------------------------------
% Introdução
% ----------------------------------------------------------

\chapter{Introdução}
\label{chap:Introdução}

Este projeto propõe uma nova abordagem baseada em sistemas neuro-fuzzy evolutivos para detecção e prevenção de fraudes em operações do setor agropecuário, utilizando dados reais de Guias de Trânsito Animal (GTA). A proposta explora a integração entre redes neurais artificiais e sistemas fuzzy evolutivos, permitindo que o modelo aprenda automaticamente funções de pertinência e regras de decisão a partir dos dados, mesmo diante de incertezas, ruídos e padrões em constante transformação. Essa arquitetura neuro-fuzzy evolutiva é especialmente adequada para ambientes dinâmicos, pois combina a capacidade de adaptação e generalização das redes neurais com a interpretabilidade e flexibilidade dos sistemas fuzzy, tornando-se robusta frente à evolução dos padrões de fraude.

O contexto do setor agropecuário brasileiro, marcado por grandes volumes de dados e vulnerabilidade a práticas fraudulentas, exige soluções inteligentes e adaptativas. A utilização de dados reais de GTAs, que registram movimentações de animais e são fundamentais para o controle sanitário e fiscal, possibilita a criação de um sistema capaz de identificar irregularidades e prevenir fraudes, como as observadas em operações citadas nesse como CIRA, Rei do Gado e Cacus.

O trabalho está estruturado para abordar desde a fundamentação teórica dos sistemas neuro-fuzzy evolutivos, passando pela análise e preparação dos dados, até a implementação, validação e discussão dos resultados obtidos. Espera-se contribuir para o avanço científico e prático na detecção de fraudes em ambientes complexos e dinâmicos, fornecendo suporte efetivo à fiscalização, à tomada de decisão e à integridade do setor agropecuário.

Os órgãos de defesa agropecuária são responsáveis pela coordenação das atividades de defesa agropecuária, com foco na segurança alimentar, saúde pública e qualidade dos produtos agropecuários. Essas instituições desempenham funções como fiscalização sanitária, controle de doenças e pragas, rastreabilidade de produtos e apoio técnico ao setor agropecuário. Muitas dessas organizações utilizam sistemas de gerenciamento do setor agropecuário que facilitam o controle sanitário, rastreabilidade e monitoramento de doenças, promovendo maior transparência e eficiência na fiscalização, além de fortalecer a competitividade do setor agropecuário. 

No sistema de gerenciamento do setor agropecuário, são coletadas diversas informações, criando um vasto banco de dados que já foi utilizado em trabalhos acadêmicos. Esses sistemas geram volumes significativos de dados, tratando-se de uma base de dados crescente que no momento conta com uma estrutura de aproximadamente 8.298.490 registros de Guias de Trânsito Animal (GTA) apenas para movimentações internas em determinadas regiões. Esses dados têm sido empregados em diversas pesquisas, como as que exploram a movimentação de bovinos, com base nas GTAs.

Nesse aspecto, o estudo de \citeonline{costa2023} analisa as redes de movimentação de bovinos entre municípios de Minas Gerais, utilizando dados da GTA coletados entre 2013 e 2023. A pesquisa identificou padrões de movimentação concentrados na região do Triângulo Mineiro e Alto Paranaíba, classificando a rede como altamente conectada, com características típicas de grafos do tipo small-world. 

Outro trabalho que pode ser citado é o artigo de \citeonline{azevedo2022}, que caracteriza a rede de movimentação de bovinos no estado de Goiás entre 2010 e 2016, utilizando dados de 4.697.239 Guias de Trânsito Animal (GTA). Os dados foram submetidos a análises estatísticas descritivas e de redes, revelando um intenso fluxo de animais, especialmente nas regiões noroeste e sul do estado, com destinos voltados para engorda, recria e abate. 

Esses dois estudos demonstram como a aplicação de análises de redes e outras ferramentas tecnológicas pode contribuir para melhorar a eficiência e a eficácia das estratégias de controle sanitário e epidemiológico, promovendo a sustentabilidade e a competitividade do setor agropecuário brasileiro.

O setor agropecuário brasileiro, especialmente o segmento de gado bovino, tem se mostrado vulnerável a práticas fraudulentas, como evidenciado em diversas operações de combate à sonegação fiscal e fraudes no comércio de gado. Podemos citar alguns exemplos de deflagrações fraudulentas encontradas na íntegra, como a:

\begin{itemize}
\item Operação do CIRA, deflagrada pelo Ministério Público de Minas Gerais em 2022, que, ao longo de 11 meses, emitiu 21 notas fiscais Guias de Trânsito Animal em nome do produtor rural sem que ele soubesse, totalizando uma quantia de R\$ 1.670.451,50. 
\item A Operação Rei do Gado, conduzida pela Receita Federal, investigou fraudes fiscais de R\$ 1,4 bilhão em transações ilegais envolvendo a venda de bovinos \cite{cnn2023}. 
\item A Operação Boi na Linha \cite{g12016} também destacou as dificuldades de fiscalização no setor, com a prisão de envolvidos em fraudes envolvendo a venda e o transporte ilegal de gado. 
\item Por fim, a recente Operação Cacus, realizada pela Polícia Civil de Minas Gerais, investigou servidores estaduais suspeitos de estarem envolvidos em um esquema de movimentação ilegal de 83 mil cabeças de gado \cite{g12024}.
\end{itemize}

Essas operações exemplificam a necessidade urgente de implementação de tecnologias inteligentes para monitorar, identificar e prevenir fraudes nesse setor. A utilização de modelos baseados em sistemas inteligentes, como o que se propõe neste estudo, poderá contribuir significativamente para reduzir as irregularidades, melhorar a fiscalização e garantir maior integridade ao sistema de emissão de documentos e transações de bovinos.

Para tanto, podemos também exemplificar alguns trabalhos existentes que investigam fraudes em diversos setores, utilizando diferentes tipos de ferramentas e propondo novos paradigmas para predizer possíveis fraudes no sistema. Alguns desses trabalhos são citados a seguir. 

O estudo de \citeonline{bao2019} propõe um modelo de previsão de fraudes contábeis utilizando técnicas de sistemas inteligentes, destacando a abordagem de ensemble learning. Em vez de utilizar razões financeiras tradicionais, os autores empregam dados financeiros brutos como entrada para o modelo que também introduziu métricas de avaliação, como o Normalized Discounted Cumulative Gain (NDCG), para priorizar casos de fraude com maior probabilidade de ocorrência, tornando-o prático para reguladores e auditores com recursos limitados. 

O trabalho de \citeonline{singh2019} apresenta um modelo de detecção de fraudes em cartões de crédito baseado em Modelos Ocultos de Markov (HMM). A proposta envolve treinar o HMM com o comportamento normal dos usuários em transações financeiras, possibilitando identificar fraudes por meio de desvios nos padrões de gasto. O modelo é dividido em duas fases: treinamento, onde os dados históricos do usuário são utilizados para criar um perfil transacional, e detecção, que analisa novas transações em tempo real. 

O artigo de \citeonline{coelho2006} apresenta o FControl®, um sistema computacional híbrido baseado em inteligência computacional, projetado para detecção de fraudes em transações de comércio eletrônico com cartões de crédito. O FControl® utiliza uma abordagem híbrida de inteligência computacional baseada em sistema neuro-nebuloso (neuro-fuzzy) e otimização por estratégia evolutiva com mecanismos de auto-adaptação de parâmetros. A prevenção de fraude em cartão de crédito representa uma importante aplicação comercial para métodos de previsão e inteligência computacional, combinando sistemas neuro-nebulosos, redes neurais artificiais, computação evolutiva e estratégias de otimização. O sistema se destaca pela rapidez e eficiência na detecção de fraudes, demonstrando a eficácia de abordagens híbridas que integram diferentes paradigmas de sistemas inteligentes para problemas de classificação e detecção de anomalias em tempo real. 

A pesquisa de \citeonline{mallela2024} explora o uso de técnicas de Machine Learning (ML) para aprimorar os mecanismos de detecção de fraudes em instituições financeiras. A pesquisa avalia modelos supervisionados, não supervisionados e de aprendizado por reforço, destacando a importância da engenharia de características e do pré-processamento de dados para a robustez dos modelos. 

No trabalho apresentado por \citeonline{hooda2018}, que apresenta uma abordagem para classificação de empresas fraudulentas utilizando algoritmos de sistemas inteligentes em auditorias externas. O estudo utilizou dados de 777 empresas de 14 setores diferentes, aplicando o algoritmo Particle Swarm Optimization (PSO) para seleção de características relevantes e comparando dez modelos de classificação. 

\citeonline{zhou2023} apresentaram uma abordagem inovadora para detecção de fraudes em demonstrações financeiras utilizando um grafo de conhecimento de duas camadas (FSFD-TLKG). A camada semântica do grafo modela os relacionamentos de subordinação entre variáveis financeiras, enquanto a camada sintática explora as relações de articulação inerentes às demonstrações financeiras. 

Esses estudos demonstram novos e combinação de diferentes métodos, como aprendizado supervisionado, não supervisionado, aprendizado por reforço, IA e grafos de conhecimento, permite a criação de sistemas robustos, eficientes e dinâmicos. Além disso, a aplicação dessas tecnologias em setores como auditoria, comércio eletrônico e finanças pode melhorar substancialmente a capacidade de identificar fraudes e reduzir os prejuízos associados. O desenvolvimento de soluções preditivas e automatizadas para detecção de fraudes está se tornando cada vez mais relevante, alinhando-se com as necessidades de eficiência e segurança nos sistemas modernos.

No contexto específico de sistemas evolutivos para detecção de fraudes, a incorporação de métodos de pertinência fuzzy emerge como uma abordagem promissora. O artigo de \citeonline{coelho2006}, por exemplo, apresenta o FControl®, um sistema híbrido neuro-fuzzy para detecção de fraudes. Os sistemas fuzzy evolutivos combinam a capacidade de lidar com incertezas e imprecisões inerentes aos dados agropecuários, por meio de funções de pertinência, com a adaptabilidade estrutural dos algoritmos evolutivos. Essa combinação permite que o sistema modele adequadamente as nuances dos padrões fraudulentos, que frequentemente apresentam características ambíguas e limites não bem definidos.

Os sistemas neuro-fuzzy evolutivos, pioneiramente desenvolvidos por Kasabov e Song \cite{kasabov-2002} com o DENFIS (Dynamic Evolving Neural-Fuzzy Inference System), representam uma evolução natural desta abordagem. Angelov e Filev \cite{angelov-2004b} expandiram esta linha de pesquisa com métodos de identificação online de modelos Takagi-Sugeno, enquanto Lughofer \cite{Lughofer2011} consolidou as metodologias e conceitos avançados desses sistemas evolutivos. Mais recentemente, trabalhos como o de Skrjanc et al. \cite{Skrjanc2019a} demonstraram a eficácia de classificadores evolutivos em ambientes de streaming de dados, características essenciais para aplicações em tempo real.

Portanto, este trabalho de pesquisa propõe o desenvolvimento de um modelo inteligente para detecção e prevenção de fraudes em sistemas de gestão, com validação por meio de dados gerados por sistemas de gerenciamento do setor agropecuário, a partir de um estudo de caso real. A pesquisa atende a uma necessidade do setor de criar métodos para mitigar fraudes no sistema. A proposta se destaca por sua abordagem específica, adaptada às particularidades do sistema agropecuário em questão, a partir de análise dos dados coletados ao longo dos anos fornecerá insights valiosos para a criação de um sistema inteligente capaz de detectar irregularidades e prevenir fraudes, como as observadas nas operações CIRA, Rei do Gado e Cacus.

O projeto está organizado da seguinte maneira: na Seção \ref{chap:Problema} é apresentada uma caracterização do problema; na Seção \ref{chap:objetivos} são definidos os objetivos gerais e específicos; na Seção \ref{chap:metodologia} descreve-se a metodologia que será utilizada no desenvolvimento da pesquisa; e na Seção \ref{chap:cronograma} são apresentados o cronograma de execução e a formação acadêmica realizada.

%----------------------------------------------------------
 
%----------------------------------------------------------
% Descrição do Problema
%----------------------------------------------------------
\chapter{Caracterização do Problema}
\label{chap:Problema}

Sistemas neuro-fuzzy evolutivos têm se destacado como uma abordagem inovadora e promissora para o tratamento de problemas complexos e dinâmicos, como a detecção de fraudes em sistemas de gestão agropecuária. Ao combinar a capacidade de adaptação dos sistemas evolutivos com o poder de representação e interpretação dos sistemas fuzzy e a aprendizagem automática das redes neurais, esses modelos são capazes de lidar com grandes volumes de dados, incertezas, ambiguidades e padrões em constante transformação. No contexto do setor agropecuário brasileiro, caracterizado por operações descentralizadas, diversidade de atores e alta variabilidade temporal, a aplicação de sistemas neuro-fuzzy evolutivos permite identificar padrões irregulares e antecipar práticas fraudulentas de forma adaptativa e eficiente. Dessa forma, o desenvolvimento e a implementação de métodos baseados em neuro-fuzzy evolutivo tornam-se essenciais para garantir a integridade, a transparência e a segurança dos processos de emissão de Guias de Trânsito Animal (GTAs) e demais operações do setor, promovendo maior confiança e sustentabilidade para toda a cadeia produtiva.

\section{Contexto e Justificativa}

Os sistemas inteligentes evolutivos têm emergido como uma abordagem promissora para problemas que envolvem fluxos contínuos de dados em ambientes dinâmicos e não-estacionários. Inicialmente tratados por Kasabov e Song \cite{kasabov-2002} e Angelov e Filev \cite{angelov-2004b}, esses sistemas possuem a característica fundamental de adaptar simultaneamente sua estrutura e parâmetros à medida que novas amostras são disponibilizadas. No contexto de detecção de fraudes, essa capacidade adaptativa é especialmente relevante, uma vez que os padrões fraudulentos tendem a evoluir constantemente para evadir sistemas de detecção.

A incorporação de métodos de pertinência fuzzy em sistemas evolutivos oferece vantagens significativas para o tratamento de dados agropecuários, que frequentemente apresentam incertezas, imprecisões e características ambíguas. As funções de pertinência fuzzy permitem modelar adequadamente situações onde os limites entre comportamentos normais e fraudulentos não são claramente definidos, proporcionando uma representação mais natural e interpretável dos padrões identificados. Os trabalhos de Lughofer \cite{Lughofer2011} e Skrjanc et al. \cite{Skrjanc2019a} demonstram a eficácia dessa abordagem em diferentes domínios de aplicação.

O setor agropecuário brasileiro apresenta particularidades que tornam a detecção de fraudes um desafio complexo. A natureza distribuída das operações, a diversidade de atores envolvidos e a variabilidade temporal dos dados criam um ambiente propício para a evolução de esquemas fraudulentos. Operações como CIRA, Rei do Gado, Boi na Linha e Cacus evidenciaram não apenas a magnitude dos prejuízos (superiores a R\$ 1,4 bilhão), mas também a sofisticação crescente dos métodos utilizados.

\section{Problema de Pesquisa e Lacunas Identificadas}

Apesar dos avanços significativos em sistemas inteligentes evolutivos, existem lacunas importantes quando se considera sua aplicação específica para detecção de fraudes em sistemas agropecuários:

\begin{enumerate}
\item \textbf{Adaptação a Dados Esparsos:} Sistemas agropecuários frequentemente apresentam dados com alta esparsidade temporal e espacial, desafiando algoritmos evolutivos tradicionais que assumem fluxos contínuos uniformes.

\item \textbf{Interpretabilidade em Contexto Regulatório:} A necessidade de explicar decisões para auditores e reguladores requer modelos com alta interpretabilidade, característica nem sempre preservada em sistemas adaptativos complexos.

\item \textbf{Detecção de Drift Conceitual:} Fraudes evoluem rapidamente, criando conceitos drifts que sistemas existentes podem não detectar adequadamente, especialmente quando as mudanças são graduais e sutis.

\item \textbf{Integração de Métodos de Pertinência Fuzzy:} A incorporação eficiente de funções de pertinência fuzzy em sistemas evolutivos adaptativos para modelagem de incertezas e ambiguidades em dados agropecuários permanece um desafio técnico significativo.

\item \textbf{Eficiência Computacional:} A necessidade de processamento em tempo real em sistemas de grande escala demanda algoritmos computacionalmente eficientes, equilibrando adaptabilidade e performance.
\end{enumerate}

O objeto de interesse deste estudo é a criação de um modelo inteligente para a detecção e prevenção de fraudes nos sistemas de gerenciamento do setor agropecuário, com foco na emissão de GTAs e nas transações relacionadas à movimentação de gado. As fraudes identificadas, como a inserção de dados falsos ou a distorção de informações para a emissão de GTAs, representam um risco significativo para a saúde pública, comprometem a eficácia da fiscalização sanitária e afetam a competitividade do setor agropecuário. Portanto, o trabalho visa desenvolver uma solução que permita a identificação de padrões irregulares e a prevenção de fraudes no sistema.

\section{Dimensões do Problema}

Podemos citar alguns pontos que podemos analisar: 

\begin{enumerate}
\item \textbf{Análise do Sistema e suas vulnerabilidades:} O trabalho começa com uma análise detalhada dos sistemas de gerenciamento do setor agropecuário, com ênfase nas funcionalidades e na estrutura do sistema, especialmente no que se refere à emissão das GTAs e à movimentação de animais. A partir desta análise, serão identificadas as vulnerabilidades e os pontos críticos que permitem a ocorrência de fraudes.

\item \textbf{Fraudes no Setor Agropecuário:} O estudo abordará as fraudes identificadas em operações anteriores. Através dessas investigações, será possível identificar os tipos de fraudes mais recorrentes e como elas afetam o sistema de defesa agropecuária.

\item \textbf{Aplicação de Modelos Inteligentes para Detecção de Fraudes:} Um dos principais focos do estudo será a proposição de modelos de sistemas inteligentes para detectar padrões irregulares e prever fraudes no sistema. A literatura existente, que utiliza modelos como aprendizado supervisionado, não supervisionado, aprendizado por reforço e grafos de conhecimento, servirá como base para a adaptação do modelo às particularidades do setor agropecuário.

\item \textbf{Validação do Modelo com Dados Reais:} A validação do modelo inteligente será feita utilizando dados reais de sistemas de gerenciamento do setor agropecuário, coletados ao longo dos anos. O estudo contará com uma base de dados crescente que atualmente possui uma estrutura de 8.298.490 registros de GTAs, especificamente de movimentações dentro de Minas Gerais, coletados até maio de 2025. Esses dados, especialmente as informações contidas nas GTAs com variáveis como origem, destino, quantidade de animais, finalidade e meio de transporte, servirão como base para a análise dos padrões de movimentação de bovinos e para a detecção de possíveis fraudes. A validação do modelo será um passo crucial para garantir que a solução proposta seja eficaz e aplicável no contexto prático.

\item \textbf{Impacto e Benefícios da Implementação de Modelos Inteligentes:} O estudo avaliará os impactos da implementação de modelos inteligentes para a detecção de fraudes no setor agropecuário. Serão discutidos os benefícios esperados, como o aumento da transparência, a melhoria na fiscalização sanitária e o fortalecimento da competitividade do mercado agropecuário brasileiro. A utilização de tecnologias inteligentes poderá reduzir as fraudes, melhorar a precisão das informações e garantir maior confiança no sistema de defesa agropecuária.

\item \textbf{Desafios e Perspectivas Futuras:} A pesquisa também abordará os desafios relacionados à adaptação de sistemas neuro-fuzzy evolutivos aos dados reais disponíveis, como a alta dimensionalidade, a presença de ruídos, a esparsidade temporal e a heterogeneidade das informações. Entre os principais desafios estão o ajuste dinâmico das regras e funções de pertinência diante de padrões que mudam ao longo do tempo, a necessidade de mecanismos eficientes para detecção de outliers e drifts conceituais, e a integração de múltiplas fontes de dados para enriquecer o processo de aprendizagem adaptativa. Além disso, o trabalho discutirá as perspectivas futuras para o aprimoramento desses modelos, considerando o uso de dados em tempo real, a automação do ajuste de parâmetros e a incorporação de técnicas de explicabilidade para garantir transparência e confiança nos resultados gerados.
\end{enumerate}


Ressalta-se que a abordagem neuro-fuzzy evolutiva proposta neste projeto é especialmente adequada para superar desafios como a adaptação a dados esparsos, a detecção de drifts conceituais e a necessidade de interpretabilidade, aspectos críticos para a detecção de fraudes em dados reais de GTAs. Dessa forma, este trabalho visa contribuir significativamente para o enfrentamento de um problema crítico no setor agropecuário, que é a fraude nos sistemas de gerenciamento, afetando a transparência, a eficiência da fiscalização sanitária e a competitividade do mercado. Ao propor um modelo inteligente para a detecção e prevenção dessas fraudes, o estudo busca proporcionar uma solução inovadora e eficaz para instituições de defesa agropecuária, fortalecendo a integridade do sistema e promovendo a sustentabilidade do setor. A pesquisa também se alinha com as tendências tecnológicas atuais e pode servir como referência para outras instituições e contextos.



%----------------------------------------------------------
% Objetivos
%----------------------------------------------------------
\chapter{Objetivos: Geral e Específicos}
\label{chap:objetivos}


A pesquisa proposta busca abordar as fraudes identificadas em sistemas de gerenciamento, com foco na emissão de GTA e nas transações relacionadas à movimentação de bovinos. Para isso, propõe-se o desenvolvimento e adaptação de um sistema neuro-fuzzy evolutivo, capaz de tratar dados complexos, dinâmicos e incertos, ajustando-se continuamente a novos padrões e estratégias de fraude. Essa abordagem visa detectar e prevenir fraudes de forma mais eficiente, promovendo maior integridade, transparência e confiabilidade ao sistema.

\section{Objetivo Geral}

Propor, adaptar e implementar novos algoritmos baseados em sistemas inteligentes evolutivos com métodos de pertinência fuzzy para detecção e prevenção de fraudes em sistemas de gerenciamento agropecuário, com foco na adaptabilidade, autonomia e interpretabilidade dos modelos propostos.

\section{Objetivos Específicos}
\begin{enumerate}
  \item \textbf{Realizar revisão abrangente da literatura sobre:} 
  \begin{itemize}
    \item Sistemas inteligentes evolutivos e adaptativos
    \item Sistemas neuro-fuzzy evolutivos e métodos de pertinência
    \item Técnicas de detecção de anomalias e fraudes em sistemas de gestão
    \item Algoritmos de aprendizado incremental e online
    \item Metodologias de adaptação estrutural e paramétrica em tempo real
  \end{itemize}
  
  \item \textbf{Propor e adaptar novos algoritmos evolutivos:} 
  \begin{itemize}
    \item Desenvolver algoritmos com capacidade de adaptação estrutural automática
    \item Implementar mecanismos de aprendizado incremental para dados de fluxo contínuo
    \item Integrar métodos de pertinência fuzzy para modelagem de incertezas e ambiguidades
    \item Garantir interpretabilidade e explicabilidade dos modelos propostos
  \end{itemize}
  
  \item \textbf{Adaptar técnicas ao contexto agropecuário:} 
  \begin{itemize}
    \item Personalizar algoritmos para características específicas de dados agropecuários
    \item Desenvolver métodos de tratamento de dados esparsos e ruidosos
    \item Implementar técnicas de detecção de drift conceitual em ambientes dinâmicos
  \end{itemize}
  
  \item \textbf{Validar algoritmos com dados reais:} 
  \begin{itemize}
    \item Testar modelos com datasets históricos de sistemas agropecuários
    \item Alcançar métricas de desempenho superiores aos métodos estado da arte
    \item Demonstrar capacidade de adaptação em cenários de mudança conceitual
  \end{itemize}
  
  \item \textbf{Avaliar e comparar desempenho:} 
  \begin{itemize}
    \item Comparar algoritmos propostos com métodos de referência usando métricas padronizadas
    \item Realizar análises estatísticas de significância dos resultados
    \item Avaliar eficiência computacional e escalabilidade dos métodos
  \end{itemize}
  
  \item \textbf{Publicar contribuições científicas:} 
  \begin{itemize}
    \item Produzir artigos em periódicos de alto impacto na área
    \item Apresentar resultados em conferências nacionais e internacionais
    \item Disponibilizar código e datasets para reprodutibilidade científica
  \end{itemize}
\end{enumerate}

\section{Resultados Esperados}

\begin{enumerate}
  \item Modelo inteligente para detecção de fraudes em sistemas de gerenciamento do setor agropecuário.
  \item Base de dados organizada e anonimizada para estudos futuros, com uma estrutura crescente de 8.298.490 registros de GTAs validados e estruturados.
  \item Relatórios sobre padrões fraudulentos e propostas de mitigação baseados em análise de dados reais.
  \item Publicações científicas em detecção de fraudes no setor agropecuário com validação experimental robusta.
  \item Propostas de melhorias que aumentem a eficiência dos sistemas de gerenciamento do setor agropecuário.
  \item Impacto positivo no setor, com redução de fraudes, fortalecimento da segurança alimentar e impactos na arrecadação.
  \item Algoritmos validados com base de dados crescente de 8.298.490 registros estruturados demonstrando aplicabilidade prática.
\end{enumerate}

%----------------------------------------------------------
% Metodologia
%----------------------------------------------------------

Sistemas neuro-fuzzy evolutivos representam o estado da arte para o tratamento de problemas complexos e dinâmicos, como a detecção de fraudes em grandes bases de dados. Ao combinar a capacidade adaptativa dos sistemas evolutivos, a modelagem de incertezas dos sistemas fuzzy e o poder de aprendizagem das redes neurais, esses modelos conseguem ajustar suas regras, funções de pertinência e parâmetros de decisão de forma contínua, acompanhando a evolução dos padrões de fraude. Essa abordagem permite alta sensibilidade e precisão mesmo diante de estratégias fraudulentas inéditas, tornando-se uma solução robusta, flexível e interpretável para ambientes adversariais e em constante transformação.

\section{Etapas da Metodologia}

\begin{enumerate}
\item \textbf{Revisão Bibliográfica e Fundamentação Teórica}
\begin{itemize}
\item Realizar levantamento sistemático da literatura sobre sistemas inteligentes evolutivos, incluindo trabalhos pioneiros e desenvolvimentos recentes
\item Analisar metodologias de detecção de anomalias e fraudes em sistemas complexos
\item Estudar algoritmos de aprendizado incremental e online, com foco em adaptação estrutural e paramétrica
\item Investigar técnicas de tratamento de drift conceitual e adaptação a ambientes não-estacionários
\item Identificar lacunas teóricas e práticas que possam ser exploradas na proposição de novos algoritmos
\end{itemize}

\item \textbf{Proposição e Implementação de Algoritmos}
\begin{itemize}
\item Propor novos algoritmos baseados em sistemas fuzzy evolutivos adaptados ao contexto de detecção de fraudes
\item Desenvolver mecanismos de adaptação estrutural automática (criação, fusão e remoção de regras/grupos)
\item Implementar técnicas de aprendizado incremental para processamento de dados em tempo real
\item Integrar métodos de pertinência fuzzy para tratamento de incertezas e modelagem de padrões ambíguos
\item Garantir interpretabilidade através de representação linguística e explicabilidade das decisões
\item Implementar algoritmos em linguagem Python, priorizando código modular, reutilizável e documentado
\item Desenvolver dashboards interativos para visualização dos resultados, utilizando bibliotecas Python como Dash ou Streamlit, de modo a facilitar a análise exploratória, a interpretação dos dados e a comunicação dos achados
\item Implementar um sistema web para disponibilização dos resultados, relatórios e ferramentas de análise, permitindo acesso remoto, interação com os dados e compartilhamento dos achados com diferentes públicos de interesse
\end{itemize}

\item \textbf{Validação Experimental e Análise Comparativa}
\begin{itemize}
\item Aplicar algoritmos propostos em datasets reais de sistemas agropecuários
\item Comparar desempenho com métodos estado da arte usando métricas padronizadas
\item Realizar análises estatísticas de significância dos resultados
\item Avaliar capacidade de adaptação em cenários de mudança conceitual
\item Analisar eficiência computacional e escalabilidade dos métodos propostos
\end{itemize}
\end{enumerate}


\section{Cronograma de Execução}

O planejamento temporal da pesquisa foi elaborado de modo a garantir a execução ordenada e eficiente de todas as etapas previstas, desde a fundamentação teórica até a validação dos resultados. O cronograma detalhado, incluindo as atividades de pesquisa, desenvolvimento, análise e as disciplinas obrigatórias do doutorado, encontra-se apresentado no Capítulo~\ref{chap:cronograma}. Ressalta-se que o acompanhamento das disciplinas cursadas, a cursar e aproveitáveis está centralizado naquele capítulo, evitando redundâncias neste.


\section{Aspectos Éticos}

Todos os procedimentos da pesquisa respeitarão rigorosamente os princípios éticos aplicáveis ao uso de dados sensíveis, em conformidade com a legislação vigente. Serão adotadas as seguintes diretrizes:
\begin{itemize}
    \item Anonimização e proteção integral dos dados utilizados;
    \item Observância à Lei Geral de Proteção de Dados (LGPD);
    \item Obtenção de autorizações institucionais formais para uso dos dados;
    \item Garantia de que nenhuma informação permita a identificação de indivíduos, propriedades ou organizações.
\end{itemize}


\section{Métricas de Avaliação}

Para garantir a robustez e a comparabilidade dos resultados, a avaliação dos algoritmos propostos será realizada por meio de métricas consagradas na literatura de detecção de fraudes e sistemas evolutivos. Serão consideradas métricas de precisão, classificação e testes estatísticos:


\subsection{Métricas de Precisão}
As métricas de precisão são utilizadas para avaliar o desempenho dos modelos em tarefas de previsão de valores contínuos:
\begin{itemize}
    \item \textbf{Raiz do Erro Quadrático Médio (RMSE):} Mede o erro médio das previsões realizadas pelo modelo em relação aos valores reais. Quanto menor o RMSE, melhor o desempenho do modelo, pois indica que as previsões estão mais próximas dos valores observados.
    \[RMSE = \sqrt{\frac{\sum_{t=1}^{N}(y_t - \hat{y}_t)^2}{N}}\]
    Onde $y_t$ é o valor real, $\hat{y}_t$ é o valor previsto e $N$ é o número total de observações.
    \item \textbf{Índice de Erro Não Dimensional (NDEI):} É a razão entre o RMSE e o desvio padrão dos valores previstos. Permite comparar o erro relativo ao grau de dispersão dos dados, tornando a métrica adimensional e facilitando a comparação entre diferentes conjuntos de dados.
    \[NDEI = \frac{RMSE}{\sigma(\hat{y}_t)}\]
    Onde $\sigma(\hat{y}_t)$ é o desvio padrão das previsões.
    \item \textbf{Soma Residual dos Quadrados (RSS):} Representa a soma dos quadrados das diferenças entre os valores reais e previstos. É utilizada para avaliar o ajuste do modelo: quanto menor o RSS, melhor o ajuste.
    \[RSS = \sum_{t=1}^{N}(y_t - \hat{y}_t)^2\]
\end{itemize}


\subsection{Métricas de Classificação}
As métricas de classificação são empregadas para avaliar o desempenho dos modelos em tarefas de detecção de fraudes (classificação binária):
\begin{itemize}
    \item \textbf{Precisão (Precision):} Indica a proporção de casos classificados como positivos que realmente são positivos. É útil para avaliar o desempenho do modelo em situações onde o custo de falsos positivos é alto.
    \[\text{Precisão} = \frac{TP}{TP + FP}\]
    Onde $TP$ são os verdadeiros positivos e $FP$ os falsos positivos.
    \item \textbf{Recall (Sensibilidade):} Mede a capacidade do modelo de identificar corretamente todos os casos positivos. É importante em contextos onde a detecção de todos os casos reais é prioritária.
    \[\text{Recall} = \frac{TP}{TP + FN}\]
    Onde $FN$ são os falsos negativos.
    \item \textbf{F1-score:} É a média harmônica entre precisão e recall, fornecendo uma medida única que equilibra ambos os aspectos. É especialmente útil quando há desequilíbrio entre as classes.
    \[F1 = 2 \cdot \frac{\text{Precisão} \cdot \text{Recall}}{\text{Precisão} + \text{Recall}}\]
    \item \textbf{Acurácia:} Representa a proporção de predições corretas (tanto positivas quanto negativas) em relação ao total de casos avaliados.
    \[\text{Acurácia} = \frac{TP + TN}{TP + TN + FP + FN}\]
    Onde $TN$ são os verdadeiros negativos.
    \item \textbf{Área sob a Curva ROC (AUC-ROC):} Mede a capacidade do classificador em distinguir entre as classes. Quanto mais próximo de 1, melhor o desempenho do modelo. Um valor de 0,5 indica desempenho aleatório.
\end{itemize}

\subsection{Testes Estatísticos}
Os testes estatísticos são fundamentais para validar se as diferenças observadas entre os modelos são estatisticamente significativas e não fruto do acaso. No contexto deste projeto, será utilizado o seguinte teste:
\begin{itemize}
    \item \textbf{Teste Morgan-Granger-Newbold (MGN):} Utilizado para comparar o desempenho de dois modelos preditivos, avaliando se a diferença entre os erros (resíduos) de ambos é estatisticamente significativa. O teste calcula um valor baseado na correlação entre os resíduos dos modelos, permitindo verificar se um modelo apresenta desempenho superior ao outro de forma consistente.
    \[MGN = \hat{\rho}_e \sqrt{\frac{N-1}{1-\hat{\rho}_e^2}}\]
    Onde $\hat{\rho}_e$ é a correlação entre os resíduos dos modelos comparados e $N$ é o número de observações. Valores elevados de MGN indicam maior diferença estatística entre os modelos.
\end{itemize}

\section{Base de Dados}

O estudo utilizará dados reais do sistema de gerenciamento agropecuário, especificamente registros de Guias de Trânsito Animal (GTA) coletados ao longo dos anos. A base de dados principal possui as seguintes características:

\subsection{Estrutura dos Dados}
O dataset contém as seguintes variáveis principais:

\begin{table}[h]
\centering
\caption{Estrutura da Base de Dados de GTAs}
\begin{tabular}{|l|p{8cm}|}
\hline
\textbf{Variável} & \textbf{Descrição} \\
\hline
nr\_gta & Número identificador único da Guia de Trânsito Animal \\
\hline
nr\_serie & Número de série da GTA \\
\hline
id\_gta & Identificador interno da GTA no sistema \\
\hline
dt\_emissao\_gta & Data de emissão da guia \\
\hline
cod\_produtor\_origem & Código do produtor de origem dos animais \\
\hline
cod\_produtor\_destino & Código do produtor de destino dos animais \\
\hline
id\_tipo\_finalidade\_gta & Identificador do tipo de finalidade da movimentação \\
\hline
ds\_meio\_transporte & Descrição do meio de transporte utilizado \\
\hline
cod\_municipio\_origem & Código do município de origem \\
\hline
nome\_municipio\_origem & Nome do município de origem \\
\hline
estado\_origem & Estado de origem (MG) \\
\hline
cod\_municipio\_destino & Código do município de destino \\
\hline
nome\_municipio\_destino & Nome do município de destino \\
\hline
estado\_destino & Estado de destino \\
\hline
qtd & Quantidade de animais transportados \\
\hline
\end{tabular}
\end{table}

\subsection{Potencial para Detecção de Fraudes}
Esta base de dados apresenta características ideais para o desenvolvimento e validação de algoritmos de detecção de fraudes:

\begin{itemize}
\item \textbf{Volume significativo:} Com mais de 8 milhões de registros, permite análises estatísticas robustas
\item \textbf{Diversidade temporal:} Dados históricos possibilitam análise de padrões evolutivos
\item \textbf{Informações relacionais:} Conexões entre produtores, municípios e finalidades permitem análise de redes
\item \textbf{Variáveis de controle:} Dados de quantidade, datas e localizações facilitam identificação de anomalias
\end{itemize}


\subsection{Tratamento e Preparação dos Dados}
Os dados serão submetidos aos seguintes processos:
\begin{enumerate}
    \item \textbf{Limpeza e validação:} Identificação e tratamento de valores ausentes, duplicatas e inconsistências
    \item \textbf{Anonimização:} Remoção ou codificação de informações que possam identificar produtores específicos
    \item \textbf{Engenharia de características:} Criação de variáveis derivadas para análise de padrões (frequência de movimentação, distâncias, sazonalidade)
    \item \textbf{Normalização:} Padronização de escalas para algoritmos de aprendizado de máquina
    \item \textbf{Detecção de outliers:} Identificação preliminar de transações atípicas
\end{enumerate}

\subsection{Proposta de Sistema Inteligente Fuzzy Evolutivo}

Com base na base de dados de GTAs, propõe-se o desenvolvimento de um sistema inteligente fuzzy evolutivo para detecção de fraudes em movimentações agropecuárias. O sistema seguirá as seguintes etapas:
\begin{enumerate}
    \item \textbf{Pré-processamento dos Dados:} Limpeza, anonimização e normalização das variáveis, além da criação de atributos derivados, como frequência de movimentação, distâncias e sazonalidade.
    \item \textbf{Modelagem Fuzzy Evolutiva:} Definição de variáveis linguísticas (ex: quantidade baixa/média/alta, distância curta/longa, frequência incomum), construção de funções de pertinência para representar incertezas e inicialização de regras fuzzy para padrões normais e suspeitos.
    \item \textbf{Evolução Online:} O sistema ajusta automaticamente as funções de pertinência e as regras fuzzy conforme novos dados são inseridos, utilizando algoritmos evolutivos para criação, fusão e remoção de regras. O aprendizado incremental permite adaptação contínua a novos comportamentos e fraudes.
    \item \textbf{Detecção de Fraudes:} Para cada nova GTA registrada, o sistema calcula o grau de pertinência a padrões normais e suspeitos, sinalizando casos anômalos para análise detalhada.
    \item \textbf{Visualização e Interação:} Dashboards interativos apresentam estatísticas, alertas e padrões detectados, facilitando a tomada de decisão por gestores e auditores.
\end{enumerate}

Essa abordagem permite identificar fraudes mesmo em cenários dinâmicos e incertos, aproveitando o potencial adaptativo dos sistemas fuzzy evolutivos e a riqueza de informações do banco de dados agropecuário.

Portanto, a metodologia proposta permite alcançar os objetivos gerais e específicos do projeto. Cada etapa está alinhada com os objetivos previamente definidos, desde a análise do problema e coleta de dados até a validação do modelo e proposição de melhorias. Isso garante que a pesquisa seja conduzida de forma sistemática, transparente e orientada a resultados teóricos, práticos e aplicáveis.


%----------------------------------------------------------
% Cronograma e Formação Acadêmica
%----------------------------------------------------------
\chapter{Cronograma e Formação Acadêmica}
\label{chap:cronograma}

O presente capítulo apresenta o planejamento temporal das atividades do doutorado, bem como a organização das disciplinas obrigatórias e o aproveitamento de créditos já cursados em outros programas.

\section{Cronograma de Execução}

Esta seção detalha o cronograma a ser seguido no Programa de Pós-Graduação em Modelagem Matemática e Computacional sob orientação do Prof. Dr. Alisson Marques da Silva no cumprimento das exigências necessárias à obtenção do título de Doutor em Modelagem Matemática e Computacional. Para tanto, na Tabela~\ref{tab:cronograma} apresenta-se o cronograma de atividades proposto para realização deste projeto, considerando o prazo máximo de 8 semestres para sua integralização.

A pesquisa será desenvolvida ao longo de 48 meses (8 semestres), dividida nas seguintes etapas:

\begin{table}[h]
\centering
\caption{Cronograma de Atividades da Pesquisa}
\label{tab:cronograma}
\begin{tabular}{|l|c|c|c|c|c|c|c|c|}
\hline
\textbf{Atividade / Semestre} & \textbf{1} & \textbf{2} & \textbf{3} & \textbf{4} & \textbf{5} & \textbf{6} & \textbf{7} & \textbf{8} \\
\hline
Integralização dos créditos & X & & & & & & & \\
\hline
Revisão Bibliográfica & X & X & X & X & X & X & & \\
\hline
Proposição de algoritmos & & X & X & X & X & & & \\
\hline
Implementação e testes & & & X & X & X & X & & \\
\hline
Validação experimental & & & & X & X & X & X & \\
\hline
Qualificação & & & & X & & & & \\
\hline
Publicações científicas & & & X & X & X & X & X & \\
\hline
Escrita e defesa da tese & & & & & & X & X & X \\
\hline
\end{tabular}
\end{table}

\subsection{Descrição das Etapas}

\begin{itemize}
\item \textbf{Semestre 1:} Integralização de créditos e início da revisão bibliográfica
\item \textbf{Semestres 2-3:} Aprofundamento da fundamentação teórica e início da proposição de algoritmos
\item \textbf{Semestres 4-5:} Desenvolvimento e implementação dos algoritmos neuro-fuzzy evolutivos
\item \textbf{Semestre 4:} Exame de qualificação
\item \textbf{Semestres 6-7:} Validação experimental e análise comparativa
\item \textbf{Semestres 6-8:} Redação da tese e preparação para defesa
\item \textbf{Semestres 3-7:} Produção contínua de publicações científicas
\end{itemize}

\section{Disciplinas do Doutorado e Planejamento}

Esta seção detalha as disciplinas obrigatórias do doutorado, sua situação atual e como elas se encaixam no planejamento acadêmico do curso.

\subsection{Disciplinas do Doutorado (PPGMMC/CEFET-MG)}

O doutorado em Modelagem Matemática e Computacional (PPGMMC/CEFET-MG) exige o cumprimento de disciplinas obrigatórias e optativas, além de atividades de pesquisa e defesa. Abaixo, estão listadas as disciplinas já cursadas (inclusive como aluno especial) e aquelas ainda pendentes:

\begin{table}[h]
\centering
\caption{Disciplinas do Doutorado – Cursadas e Pendentes}
\label{tab:disciplinas_doutorado}
\resizebox{\textwidth}{!}{%
\begin{tabular}{|l|l|l|c|c|c|}
\hline
\textbf{Disciplina} & \textbf{Professor} & \textbf{Situação} \\
\hline
Algoritmos e Estruturas de Dados & Thiago de Souza Rodrigues & Obrigatória/Cursando (Isolada) \\
\hline
Inteligência Computacional & Alisson Marques da Silva & Obrigatória/Cursando (Isolada) \\
\hline
Álgebra Linear & -- & Obrigatória/A cursar \\
\hline
Princípios de Modelagem Matemática & -- & Obrigatória/A cursar \\
\hline
Engenharia de Software & -- & Obrigatória/A cursar \\
\hline
Planejamento e Análise Estatística de Experimentos & -- & Obrigatória/A cursar \\
\hline
Desenvolvimento de Projeto de Tese I & -- & Obrigatória/A cursar \\
\hline
Desenvolvimento de Projeto de Tese II & -- & Obrigatória/A cursar \\
\hline
Desenvolvimento de Projeto de Tese III & -- & Obrigatória/A cursar \\
\hline
Elaboração de Projeto de Tese & -- & Obrigatória/A cursar \\
\hline
Defesa de Tese & -- & Obrigatória/A cursar \\
\hline
Exame de Qualificação - Doutorado & -- & Obrigatória/A cursar \\
\hline
\end{tabular}%
}
\end{table}


\section{Formação Acadêmica e Disciplinas Cursadas}

Esta seção apresenta a formação acadêmica complementar realizada em diferentes programas de pós-graduação, demonstrando o embasamento teórico multidisciplinar para o desenvolvimento desta pesquisa.\newline
\textbf{Mestrado Profissional em Modelagem Computacional e Sistemas} pela Universidade Estadual de Montes Claros (UNIMONTES) -- PPGMCS.\newline
Na UFMG, o curso de doutorado nessa área se chama \textbf{Programa de Pós-graduação em Ciência da Computação} (PPGCC), oferecido pelo Departamento de Ciência da Computação (DCC), vinculado ao Instituto de Ciências Exatas (ICEx).


\subsection{Disciplinas Cursadas no Mestrado Profissional (UNIMONTES)}

No Mestrado Profissional em Modelagem Computacional e Sistemas da Universidade Estadual de Montes Claros (UNIMONTES) -- PPGMCS, foram cursadas disciplinas fundamentais que forneceram a base teórica para o desenvolvimento de sistemas inteligentes, modelagem matemática e métodos de detecção de anomalias:

\begin{table}[ht]
\centering
\caption{Disciplinas Cursadas no Programa de Mestrado Profissional (UNIMONTES)}
\label{tab:disciplinas_mestrado}
\resizebox{\textwidth}{!}{%
\begin{tabular}{|l|l|l|c|}
\hline
\textbf{Semestre/Ano} & \textbf{Disciplina} & \textbf{Professor} & \textbf{Nota} \\
\hline
2016/2 & Métodos Matemáticos & Prof. Dr. Gustavo Foscolo de Moura Gomes & 98 \\
\hline
2016/2 & Detecção e Diagnóstico de Falhas em Sistemas Dinâmicos & Prof. Dr. Marcos Flávio Silveira Vasconcelos D'Angelo & 91 \\
\hline
2016/2 & Processos Estocásticos & Prof. Dr. Nilson Luiz Castelucio Brito & 90 \\
\hline
2017/1 & Projeto e Análise de Algoritmos & Prof. Dr. René Rodrigues Veloso & 80 \\
\hline
2017/1 & Inteligência Computacional & Prof. Dr. Marcos Flávio Silveira Vasconcelos D'Angelo & 88 \\
\hline
2017/2 & Visão Computacional & Prof. Dr. Antonio Vilson Vieira & 86 \\
\hline
\end{tabular}%
}
\end{table}

\subsection{Disciplinas Isoladas (UFMG)}

Para complementar a formação e aprofundar conhecimentos específicos relacionados ao projeto, foram cursadas disciplinas isoladas na Universidade Federal de Minas Gerais:

\begin{table}[h]
\centering
\caption{Disciplinas Isoladas Cursadas na UFMG}
\label{tab:disciplinas_ufmg}
\resizebox{\textwidth}{!}{%
\begin{tabular}{|l|l|l|c|}
\hline
\textbf{Semestre/Ano} & \textbf{Disciplina} & \textbf{Professor} & \textbf{Nota} \\
\hline
2021/2 & Mineração de Dados & Prof. Dr. Wagner Meira Júnior & 86,0 \\
\hline
2021/2 & Tópicos Especiais em Ciência da Computação & Prof. Dr. Cristiano Arbex Valle & 95,0 \\
& \textit{(Finanças Quantitativas e Gerenciamento de Risco)} & & \\
\hline
2021/1 & Tópicos Especiais em Ciência da Computação & Prof. Dr. Erickson Rangel do Nascimento & 75,0 \\
& \textit{(Visão Computacional)} & & \\
\hline
2021/1 & Tópicos Especiais em Ciência da Computação & Profa. Dra. Raquel Cardoso de Melo Minardi & 99,0 \\
& \textit{(Visualização de Dados)} & & \\
\hline
\end{tabular}%
}
\end{table}


Esta formação abrangente, aliada ao cronograma estruturado de pesquisa, estabelece as condições ideais para o desenvolvimento bem-sucedido da investigação em detecção de fraudes utilizando sistemas neuro-fuzzy evolutivos aplicados ao contexto proposto. O aproveitamento de disciplinas já cursadas permitiu reduzir o tempo dedicado à formação básica, possibilitando maior foco no desenvolvimento, implementação e validação dos algoritmos propostos. Além disso, a experiência adquirida em áreas como modelagem matemática, inteligência computacional, análise de dados e mineração de informações é fundamental para a condução de uma pesquisa inovadora, robusta e alinhada às demandas atuais do setor. Dessa forma, o planejamento acadêmico e científico apresentado neste capítulo garante a viabilidade e a relevância do projeto, contribuindo para avanços significativos na detecção de fraudes e para a formação qualificado.


%----------------------------------------------------------

% ----------------------------------------------------------
% ELEMENTOS PÓS-TEXTUAIS
% ----------------------------------------------------------
\postextual

% Referências bibliográficas
\include{./03-Elementos-pos-textuais/referencias}

% Índice remissivo
\include{./03-Elementos-pos-textuais/indiceRemissivo}

\end{document}
